% Replace pilot counts/point estimates in VI-A and VI-D; use direct-session-start-v2 model definitions.
All sessions start locally at R; latency is measured from local session activation to correction completion. Only R--B result delivery contributes classical delay.

The expanded finite-batch study uses 32 independent traffic-phase seeds,
two quantum seeds, three offered loads, and two request intervals, yielding
384 paired cases and 1,536 transactions per model. Calibration remains
fixed from two disjoint traffic seeds. We report 95\% percentile-bootstrap
intervals from 10,000 resamples of traffic-seed clusters; session and quantum
repetitions are averaged within each cluster. These intervals are conditional
on the fixed calibration and the specified periodic-traffic workload family.

\begin{table*}[t]
\centering
\caption{Latency MAE (ms), mean [95\% CI], at a 0.8-ms request interval.}
\begin{tabular}{lccc}
\hline
Model & Load 0 & Load 0.35 & Load 0.70 \\
\hline
Fixed-Dc & 0.000 [0.000, 0.000] & 0.174 [0.144, 0.206] & 0.242 [0.227, 0.258] \\
No-Dq-R & 1.200 [1.200, 1.200] & 1.188 [1.185, 1.191] & 1.188 [1.174, 1.208] \\
Decoupled & 0.000 [0.000, 0.000] & 0.117 [0.073, 0.165] & 0.223 [0.189, 0.260] \\
\hline
\end{tabular}
\end{table*}

At a 0.8-ms request interval, the expected-fidelity biases of No-Dq-R are 0.0435 [0.0435, 0.0435], 0.0423 [0.0419, 0.0426], 0.0409 [0.0403, 0.0415], respectively. These are probability-weighted branch expectations, not single observed-branch fidelities.

At a 2-ms interval, the Fixed-Dc latency MAEs are 0.184 [0.151, 0.216] and 0.240 [0.227, 0.253] ms at loads 0.35 and 0.70. No-Dq-R and Decoupled have zero latency error in these noncontending finite workloads.

With the unchanged $F_{\min}=0.5$ and 5-ms deadline measured from local session start, No-Dq-R has false service-feasible fractions 0.2500 [0.2500, 0.2500], 0.1875 [0.1484, 0.2266], 0.0938 [0.0547, 0.1328] at a 0.8-ms interval. These are finite-workload, paired sampled classifications; they are not hardware success probabilities. The zero-load traffic phase is inactive, so its degenerate interval does not establish population certainty.

\begin{table*}[t]
\centering
\caption{Service-feasibility misclassification fractions, mean [95\% CI]. Latency is measured from local session start.}
\begin{tabular}{lllcc}
\hline
Interval (ms) & Load & Model & False feasible & False infeasible \\
\hline
0.8 & 0.0 & Fixed-Dc & 0.0000 [0.0000, 0.0000] & 0.0000 [0.0000, 0.0000] \\
0.8 & 0.0 & No-Dq-R & 0.2500 [0.2500, 0.2500] & 0.0000 [0.0000, 0.0000] \\
2.0 & 0.0 & Fixed-Dc & 0.0000 [0.0000, 0.0000] & 0.0000 [0.0000, 0.0000] \\
2.0 & 0.0 & No-Dq-R & 0.0000 [0.0000, 0.0000] & 0.0000 [0.0000, 0.0000] \\
0.8 & 0.35 & Fixed-Dc & 0.0000 [0.0000, 0.0000] & 0.0000 [0.0000, 0.0000] \\
0.8 & 0.35 & No-Dq-R & 0.1875 [0.1484, 0.2266] & 0.0000 [0.0000, 0.0000] \\
2.0 & 0.35 & Fixed-Dc & 0.0000 [0.0000, 0.0000] & 0.0000 [0.0000, 0.0000] \\
2.0 & 0.35 & No-Dq-R & 0.0000 [0.0000, 0.0000] & 0.0000 [0.0000, 0.0000] \\
0.8 & 0.7 & Fixed-Dc & 0.0000 [0.0000, 0.0000] & 0.0000 [0.0000, 0.0000] \\
0.8 & 0.7 & No-Dq-R & 0.0938 [0.0547, 0.1328] & 0.0000 [0.0000, 0.0000] \\
2.0 & 0.7 & Fixed-Dc & 0.0000 [0.0000, 0.0000] & 0.0000 [0.0000, 0.0000] \\
2.0 & 0.7 & No-Dq-R & 0.0000 [0.0000, 0.0000] & 0.0000 [0.0000, 0.0000] \\
\hline
\end{tabular}
\end{table*}

Deadline-only false-feasible rates are reported separately below; they must not be interpreted as service-feasibility errors. Rates use all paired transactions in each cell as the denominator.

\begin{table*}[t]
\centering
\caption{Deadline-only false-feasible fraction, mean [95\% CI], at 0.8 ms.}
\begin{tabular}{lcc}
\hline
Model & Load 0.35 & Load 0.70 \\
\hline
Fixed-Dc & 0.0156 [0.0000, 0.0391] & 0.0234 [0.0000, 0.0547] \\
No-Dq-R & 0.0156 [0.0000, 0.0391] & 0.0234 [0.0000, 0.0547] \\
Decoupled & 0.0000 [0.0000, 0.0000] & 0.0000 [0.0000, 0.0000] \\
\hline
\end{tabular}
\end{table*}
